\chapter{Introdução}
\label{cap:introducao}

A operação segura, confiável e econômica de Sistemas Elétricos de Potência (SEP) é um dos maiores desafios da engenharia moderna. Com o crescimento contínuo da demanda por energia, a inserção massiva de fontes renováveis intermitentes e a dimensão continental do Sistema Interligado Nacional (SIN), a análise de redes elétricas tornou-se um problema de larga escala cuja resolução depende, simultaneamente, de modelos matemáticos rigorosos e de implementações computacionais eficientes. Nesse cenário, o estudo de Fluxo de Potência (FP) e a sua extensão como problema de otimização, denominada Fluxo de Potência Ótimo (FPO), são as duas ferramentas centrais sobre as quais o planejamento e a operação dos SEP são executados \cite{vonMeier2006powersystems,gloverSarma2012power}.

O FP em corrente alternada (CA) é classicamente formulado como um sistema de equações algébricas não lineares e não convexas, derivado da aplicação das Leis de Kirchhoff na rede elétrica. O FPO acrescenta a esse sistema uma função objetivo, tipicamente o custo de geração ou as perdas ativas totais, e um conjunto de restrições de desigualdade que codificam os limites operacionais dos equipamentos. Em sua forma natural, o FPO pertence à classe de problemas de programação não linear (NLP) não convexos, cuja resolução em redes de grande porte exige métodos numéricos avançados, formulações cuidadosas e \textit{solvers} robustos \cite{molzahn2019survey}.

\section{Motivação}
\label{sec:motivacao}

A análise de SEP no contexto brasileiro impõe particularidades que vão além do FPO acadêmico. A operação do SIN é fortemente dependente de \textbf{ações de controle}, tais como limites de geração reativa em barras PV (QLIM), limites de magnitude de tensão (VLIM), bancos \textit{shunt} chaveados (CSCA), \textit{taps} ajustáveis de transformadores em carga (CTAP) e, em condições de emergência, esquemas hierárquicos de corte/alívio de carga (rotulados neste trabalho como DERA). Essas ações precisam ser representadas de forma explícita para que os pontos de operação calculados em ambiente computacional reproduzam um comportamento mais próximo das condições reais de operação do sistema. Programas consolidados na indústria nacional, com destaque para o ANAREDE do Centro de Pesquisas de Energia Elétrica (CEPEL)~\cite{cepel2023anarede}, incorporam esses controles por meio de heurísticas e laços iterativos específicos, codificados em uma cultura de modelagem que está documentada principalmente nos manuais e em arquivos de dados no formato PWF.

Por outro lado, o ecossistema acadêmico de pesquisa em otimização de SEP tem se consolidado em torno de plataformas \textit{open-source} de propósito geral, destacando-se o \textit{PowerModels.jl}~\cite{coffrin2018powermodels}, desenvolvido em linguagem Julia e amplamente utilizado na pesquisa internacional para estudos de Fluxo de Potência Ótimo (FPO). Embora essa ferramenta ofereça formulações robustas para o FPO AC e para suas relaxações DC, SOC e SDP, ela não incorpora, de forma nativa, os mecanismos de controle amplamente empregados nos estudos de operação realizados no SIN, disponíveis no ANAREDE. Surge, assim, uma lacuna metodológica relevante: \emph{como reproduzir, em uma ferramenta moderna e \textit{open-source} da linguagem Julia, como o JuMP, o conjunto de ações de controle que tornam o FPO compatível com a prática de operação do SIN?}

Essa lacuna motivou o desenvolvimento do pacote \textit{ControlPowerFlow.jl}, uma implementação modular dessas ações de controle em Julia, disponível publicamente no GitHub \cite{chavarry2024}. Segundo os autores, porém, o protótipo do \textit{ControlPowerFlow.jl} foi descontinuado e está incompleto, o que inviabilizou seu uso como referência de comparação direta no presente estudo.

Por conta disso, este trabalho parte de uma abordagem independente: são desenvolvidas desde a concepção, em \textit{JuMP}, formulações de FP e FPO que incorporam, uma a uma, as mesmas ações de controle previstas no \textit{ControlPowerFlow.jl} (QLIM, VLIM, CSCA, CTAP e DERA), de forma que cada formulação possa ser comparada contra as duas referências disponíveis no \textit{PowerModels.jl} (FPO clássico e FP clássico) sobre o mesmo conjunto de casos de teste. Além de validar os resultados, essa abordagem permite documentar a lógica de cada ação de controle de forma explícita e reprodutível. Esses fatores, em conjunto, contribuíram para o desenvolvimento do presente trabalho.

\section{Objetivos}
\label{sec:objetivos}

\subsection{Objetivo geral}
\label{subsec:obj-geral}

Este trabalho tem como objetivo geral apresentar a implementação, o teste e a validação de uma arquitetura de progressão incremental de ações de controle de fluxo de potência ótimo em linguagem \textit{open-source} Julia, incorporando de forma explícita as principais ações de controle utilizadas pelo programa ANAREDE (QLIM, VLIM, CSCA, CTAP e DERA).

\subsection{Objetivos específicos}
\label{subsec:obj-especificos}

\begin{alineascomponto}
    \item compreender e documentar a formulação matemática do FP e do FPO em coordenadas polares, bem como as principais relaxações disponíveis na literatura;
    \item utilizar a biblioteca \textit{PWF.jl}~\cite{pwfjl} para a leitura de arquivos no formato ANAREDE, com extração das estruturas de controle (registros DOPC e DLIN);
    \item implementar duas formulações de referência baseadas em \textit{PowerModels.jl}, uma de FPO completo (\textit{solve\_ac\_opf}) e uma de FP clássico (\textit{solve\_ac\_pf}), para servirem como referência de comparação;
    \item implementar quatro formulações desenvolvidas manualmente em JuMP que adicionam, sequencialmente, as ações de controle QLIM+VLIM, CSCA, CTAP e DERA, sempre por meio de restrições do tipo \textit{soft-constraint} --- isto é, restrições flexíveis que são incorporadas à função objetivo por meio de variáveis de folga penalizadas, permitindo violações controladas para garantir a convergência numérica;
    \item submeter as seis formulações a um conjunto padronizado de vinte e sete casos de teste, abrangendo redes pedagógicas, redes de médio porte, recortes reduzidos do SIN e o cenário integral do SIN, e analisar o comportamento de convergência, o agrupamento das soluções em \textit{clusters} numericamente equivalentes e os desvios barra a barra em relação às referências.
\end{alineascomponto}

\section{Organização do trabalho}
\label{sec:organizacao}

Para uma melhor compreensão e encadeamento lógico, este documento está estruturado em cinco capítulos. O \autoref{cap:fundamentacao-teorica} caracteriza o SIN e seu arcabouço normativo, apresenta a fundamentação teórica do FP e do FPO em coordenadas polares, discute as principais relaxações da literatura, revisa as ações de controle que motivam o trabalho e encerra com uma revisão bibliográfica dos trabalhos de referência. O \autoref{chap:metodologia} detalha a metodologia adotada: a arquitetura de software em Julia, o uso do \textit{PWF.jl} para leitura de dados, a configuração do \textit{solver} Ipopt e, principalmente, a formulação matemática das seis variantes de FP/FPO implementadas. No \autoref{chap:resultados} são expostos os resultados do conjunto de casos de teste: a tabela global de convergência, o agrupamento das soluções em \textit{clusters} numericamente equivalentes e a análise de desvios barra a barra. Por fim, o \autoref{chap:conclusoes-e-trabalhos-futuros} compila as conclusões alcançadas e propõe direcionamentos para trabalhos futuros.
